\scanpage{227}

\mspage{101}

Par la condition, pour $A$ dans $\mathrm{Cat}$

\[
A\in\mathrm{Filt}^{\Pi}(\mathrm{Cat})=\mathrm{Cat}^{\Pi}\quad\Longleftrightarrow\quad\mathrm{card}\,\mathrm{Fl}(A)\leq\Pi\tag*{(8.4)}
\]

Il faut vérifier \deleted{\unsure{des}} les conditions a) b) c) de loc. cit 9.12. La condition a) est évidente, b) est contenue dans la condition plus forte de stabilité de $\mathrm{Cat}^{\Pi}$ par $\varinjlim_I$, pour $I$ une cat. telle que $\mathrm{Card}\,\mathrm{Fl}\,I\leq\Pi$ (c'est la condition d$_{\unsure{\alpha|0}}$) de ces notes, p. 45-46). Notons que $\mathrm{Cat}^{\Pi}$ est également stable par $\varprojlim$ finies, (et si $\Pi$ est de la forme $2^c$, il est stable $=$ par limites projectives $\varprojlim_I$, avec $\mathrm{card}\,\mathrm{Fl}\,I\leq c$ \ldots). La condition c) se vérifie aisément par la méthode du \S\,5, en notant que pour $X$ dans $\mathrm{Cat}$, l'ens

\[
\unsure{\mathrm{Ss}|\mathrm{Sos}}_{\Pi}(X)
\]

des sous-catégories de $X$ qui sont dans $\mathrm{Cat}^{\Pi}$ est filtrant croissant et grand devant $\Pi$, et que $X$ en est la limite inductive. D'autre part, il est clair aussi que si $X\in\mathrm{Cat}^{\Pi'}$ avec $\Pi'>\Pi$, alors \deleted{\unsure{Card}}

\[
\mathrm{Card}\,\mathrm{Ss}_{\Pi}(X)\leq{\Pi'}^{\Pi}
\]

ce qui achève d'établir la condition c).

De plus, la condition d) de p. 45 se vérifie immédiatement \deleted{pour} les $\mathrm{Cat}^{\Pi}$, et aussi pour les $\widehat A^{\Pi}$ (pour $\Pi\geq\mathrm{Card}\,\mathrm{Fl}\,A$).

Ainsi que les mono sont stables par $\varinjlim$ filtrantes (condition e) page 46).
% archive sheet number: 226
